
\chapter{Conclusion}
\begin{spacing}{1.3}

This project set out to answer a fundamental question: what does it take to transform a machine learning research prototype into a production-ready, edge-deployable system for melanoma detection? The answer, demonstrated across 12 implementation phases, is a disciplined, end-to-end MLOps pipeline that integrates data versioning, validation, training, explainability, optimisation, deployment, monitoring, and safety.

The pipeline is a functioning system whose every component has been validated against real data and measurable targets. The EfficientNet-B0 model, trained on 6,959 HAM10000 images with patient-aware splitting, achieves an ROC-AUC of 0.91 with sensitivity of 80.6\% (at clinical threshold 0.15) on a held-out test set of 1,527 images. It was exported to ONNX (FP32, 16.6~MB), served via a FastAPI application with 6 endpoints, containerised in a multi-stage Docker image, and deployed as a public Gradio demo and REST API on Hugging Face Spaces. Monitoring via Prometheus, structlog, and Evidently AI provides observability into every aspect of system behaviour. Safety guardrails -- confidence thresholds, uncertainty gating, and graceful fallback -- ensure the system degrades safely under ambiguity.

The transformation from a notebook classification demo to a production MLOps pipeline required approximately 80\% of total project effort -- consistent with the ``hidden technical debt'' thesis of Sculley et al.~\cite{sculley2015}. Model training accounts for perhaps 15\% of the codebase. The remaining 85\% is infrastructure: data pipelines, validation suites, experiment tracking, ONNX export and verification, API design, Docker multi-stage builds, monitoring instrumentation, CI/CD workflows, and safety mechanisms. This ratio is not a failure of efficiency -- it is the nature of production machine learning.

Several directions for future work emerge from this project. Multi-class classification would extend from binary to the full 7-class HAM10000 taxonomy, enabling differential diagnosis of specific lesion types. Fitzpatrick scale evaluation across the full I--VI range is essential for equitable deployment, as the HAM10000 dataset is biased toward lighter skin types. TensorRT on NVIDIA Jetson would offer superior performance through kernel auto-tuning for edge deployments on that platform. Active learning with human-in-the-loop feedback on flagged cases could create a virtuous cycle of continuous improvement. A/B testing frameworks would enable safe online evaluation of model updates before full rollout. Federated learning across multiple edge devices without centralising patient data would preserve privacy while leveraging distributed data. Multi-modal inputs incorporating patient metadata alongside image data have been shown to improve diagnostic accuracy in ISIC challenge submissions. Continuous training triggered by data drift detection would close the loop between monitoring and model improvement.

MLOps is what separates a model that works in a notebook from a system that works in a clinic. This project demonstrates that the gap is bridgeable with the right tools, practices, and mindset. The 87\% of machine learning models that never reach production do not fail because of bad algorithms. They fail because of untested data pipelines, unreproducible training procedures, unmonitored deployments, and unhandled edge cases. MLOps addresses each of these failure modes directly. This project is a testament to that approach -- and an illustration that production engineering must be treated as a first-class discipline in machine learning, not an afterthought.

\end{spacing}
